% Generated draft: requires mathematical and source review.
\documentclass[11pt]{article}
\usepackage[margin=25mm]{geometry}
\usepackage{fontspec}
\setmainfont{Latin Modern Roman}
\usepackage{amsmath,amssymb,mathtools,unicode-math}
\setmathfont{Latin Modern Math}
\usepackage{enumitem,needspace,xurl,hyperref}
\setlength{\parindent}{0pt}
\setlength{\parskip}{5pt}
\setlength{\emergencystretch}{3em}
\newcommand{\N}{\mathbb{N}}
\newcommand{\Z}{\mathbb{Z}}
\newcommand{\Q}{\mathbb{Q}}
\newcommand{\R}{\mathbb{R}}
\newcommand{\C}{\mathbb{C}}
\newcommand{\F}{\mathbb{F}}
\newcommand{\A}{\mathbb{A}}
\newcommand{\PP}{\mathbb{P}}
\newcommand{\E}{\mathbb{E}}
\newcommand{\eps}{\varepsilon}
\newcommand{\dd}{\,\mathrm{d}}
\begin{document}
\section*{UnsolvedMath ID: 30005926}
\textbf{Source problem code:} catalogue entry
\subsection{Definitions and mathematical statement}
Let \(T_{n,g}\) denote the random triangulation of size \(n\) and genus \(g\) appearing in the supplied source. For a finite connected triangulation \(T\), let \(V(T)\) be its vertex set, let \(d_T(u,v)\) be graph distance, with every edge having length one, and define
\[
\operatorname{diam}(T)=\max_{u,v\in V(T)}d_T(u,v).
\]
The supplied excerpt does not specify the probability law of \(T_{n,g}\), the meaning of size, rooting conventions, or whether loops and multiple edges are permitted. These choices remain part of the unresolved formulation.

The intended setting appears to be linear genus: an admissible sequence of integers \(g_n\) with \(g_n/n\to\theta>0\). For each fixed admissible \(\theta\), the source conjectures constants \(D_\theta\) and \(D'_\theta\) such that typical vertex distances are close to \(D_\theta\log n\), the diameter is close to \(D'_\theta\log n\), and
\[
D'_\theta=3D_\theta.
\]
Here \(\log\) may be taken as the natural logarithm, with the constants scaled accordingly for any other fixed base.

A precise candidate interpretation is the following, conditional on specifying the triangulation ensemble and its admissible parameter range. Given \(T_{n,g_n}\), choose \(U_n,V_n\) independently and uniformly from \(V(T_{n,g_n})\). Ask whether, for every admissible sequence with \(g_n/n\to\theta\), there exist deterministic constants \(D_\theta,D'_\theta\in(0,\infty)\), depending only on \(\theta\), such that
\[
\frac{d_{T_{n,g_n}}(U_n,V_n)}{\log n}\xrightarrow{\mathbb P}D_\theta,
\qquad
\frac{\operatorname{diam}(T_{n,g_n})}{\log n}\xrightarrow{\mathbb P}D'_\theta,
\qquad D'_\theta=3D_\theta.
\]
The first probability includes both the random triangulation and the sampled vertices. Uniform independent sampling and convergence in probability are explicit proposed interpretations of ``typical'' and ``close to''; neither is specified in the supplied excerpt. Consequently, these displays are a conditional formalization requiring review, rather than a uniquely determined transcription.

The adjoining source discussion gives an average-degree expression apparently intended as
\[
\frac{6n}{n+2-2g}\approx\frac{6}{1-2g/n}.
\]
For a closed orientable triangulation with \(2n\) triangular faces, Euler's formula gives this first expression exactly. This suggests, but does not establish, that the source uses \(n\) as half the number of faces. The second expression is an asymptotic reading away from a vanishing denominator. The usual interior range suggested by this expression is \(0<\theta<1/2\), but the excerpt's interval notation is corrupted, and neither the exact intended range nor endpoint claims can be recovered reliably. No distance assertion at \(\theta=0\) or \(\theta=1/2\) is made in this draft.

As supplied contextual claims, the excerpt reports a local limit \(T_\theta\), a random infinite plane triangulation whose distribution depends only on \(\theta\), with the uniform infinite planar triangulation at \(\theta=0\). It describes the latter as having average degree \(6\), compares positive genus density with the \(7\)-regular hyperbolic triangulation, and attributes exponential volume growth and positive random-walk speed to the positive-parameter limit models. It also reports enumeration estimates in any genus regime with multiplicative error \(e^{o(n)}\). The associated local-convergence interval is corrupted, and the excerpt's references [2] and [5] are unresolved. These contextual assertions do not specify the missing hypotheses of the distance conjecture and are not independently verified here.
\subsection{Short English statement}
The source conjectures that, for random triangulations whose genus grows linearly with size, distances between typical vertices and the diameter have logarithmic asymptotics, with the diameter constant three times the typical-distance constant. The triangulation law, size convention, parameter endpoints, vertex sampling, and convergence mode require clarification. Supplied curation labels the problem partially solved; that dated assessment is not independently verified.
\subsection{Sources}
\begin{itemize}
\item Supplied local UnsolvedMath registry record, permanent ID \(30005926\), title: \emph{Distance and Diameter Constants of High-Genus Triangulations}; exact source problem code \texttt{catalogue entry}. The record belongs to the Oberwolfach Reports Open Problems collection and assigns proposed year \(2024\); no individual proposer is supplied.
\item \emph{Statistical Physics and Random Surfaces} (2024), Oberwolfach Reports, pp.~1387--1444, DOI \texttt{10.4171/owr/2024/25}. Supplied report provenance: \url{https://doi.org/10.4171/owr/2024/25}. The supplied excerpt does not identify a problem-specific page, speaker, or author.
\item Supplied literature reference labeled ``Budzinski--Chapuy--Louf'': \url{https://doi.org/10.1214/24-AOP1746}. No title or precise result statement is supplied.
\item Supplied literature reference labeled ``Lions 2026'': \url{https://arxiv.org/abs/2606.27357}. The record describes this as a 2026 preprint; no title or precise result statement is supplied.
\item Literature assessment checked at \texttt{2026-08-22}: this is supplied dated curation, not an independently verified assessment. It labels the problem \texttt{partially\_solved}, reports published logarithmic distance scales and a 2026 preprint identifying the limiting constant for typical distances, and states that the diameter constant and the conjectured equality \(D'_\theta=3D_\theta\) remain unsettled. The record explicitly requests independent verification of time-sensitive status and cited claims.
\item The registry labels the statement \texttt{corrected\_verified} and says that only formula delimiters, notation, grammar, and typography were repaired, preserving mathematical content and scope. This supplied editorial label does not establish mathematical correctness. The excerpt also cites internal references [2] and [5], whose bibliographic identities are not supplied.
\end{itemize}
\end{document}
